% !PRINTED 455
\weberpage{455}

%<<CONT
und die Unveränderlichkeit durch $\chi$, dass neben diesem einen Gliede noch die zwei entsprechenden

\[
x_2^{h_1} x_3^{h_2} x_1^{h_3}, \quad x_3^{h_1} x_1^{h_2} x_2^{h_3}
\]

in der Function vorkommen, wenn nicht $h_1 = h_2 = h_3$ ist. Dazu kommt noch die Bedingung der Unveränderlichkeit durch $\omega$. Wir wollen nun sehen, wie wir diesen Forderungen genügen können, und zwar zunächst so, dass der Grad $m$ der Invariante, d.\,h. die Summe $h_1 + h_2 + h_3$, möglichst klein wird.

Die Bedingung (1) kann offenbar nicht erfüllt sein, wenn $m < 3$ ist. Ist diese Summe $= 3$, so muss $h_1 = h_2 = h_3 = 1$ sein; aber das Product $x_1 x_2 x_3$ ist offenbar nicht unverändert durch $\omega$. Der kleinste Werth, der in Betracht kommt, ist also $m = 4$, und es sind also alle nicht negativen Lösungen von

\[
\begin{aligned}
h_1 + h_2 + h_3 &= 4 \\
h_1 + 2h_2 + 4h_3 &\equiv 0\ (\mathrm{mod}\ 7)
\end{aligned}
\]
%CONT>>

aufzusuchen. Eliminiren wir $h_1$, so folgt

\[
h_2 + 3h_3 \equiv 3\ (\mathrm{mod}\ 7),
\]

und daraus ergeben sich die einzig möglichen Lösungen:

\[
\begin{aligned}
h_3 &= 0, & h_2 &= 3, & h_1 &= 1 \\
h_3 &= 1, & h_2 &= 0, & h_1 &= 3 \\
h_3 &= 3, & h_2 &= 1, & h_1 &= 0,
\end{aligned}
\]

und folglich die \gesperrt{einzige}, durch $\chi$ und $\tau$ ungeänderte Form $4^{\mathrm{ten}}$ Grades

\weq{2}{f\,(x_1, x_2, x_3) = x_1^3 x_3 + x_2^3 x_1 + x_3^3 x_2.}

Um den Einfluss der Substitution $\omega$ auf die Function $f$ zu prüfen, setzen wir

\weq{3}{
\begin{aligned}
x_1 &= \alpha y_1 + \beta y_2 + \gamma y_3 \\
x_2 &= \beta y_1 + \gamma y_2 + \alpha y_3 \\
x_3 &= \gamma y_1 + \alpha y_2 + \beta y_3,
\end{aligned}
}

worin die Coëfficienten $\alpha, \beta, \gamma$ die in \S.\,115 festgesetzte Bedeutung haben, und nehmen an, dass durch (3) die Transformation

\weq{4}{F\,(y_1, y_2, y_3) = f\,(x_1, x_2, x_3)}

geleistet werde. Wir bilden die zweiten Ableitungen von $F$ nach $y_1, y_2, y_3$ mit Rücksicht auf (3) und auf die Formeln

\[
\begin{aligned}
&\tfrac{1}{6} f''(x_1, x_1) = x_1 x_3, \quad \tfrac{1}{6} f''(x_2, x_2) = x_2 x_1, \quad \tfrac{1}{6} f''(x_3, x_3) = x_3 x_2 \\
&\tfrac{1}{3} f''(x_2, x_3) = x_3^2, \quad \tfrac{1}{3} f''(x_3, x_1) = x_1^2, \quad \tfrac{1}{3} f''(x_1, x_2) = x_2^2.
\end{aligned}
\]
